% Source: https://artofproblemsolving.com/community/c3h1199408_mathcounts_2010_sprint_30
% Stable identity: aops-topic-1199408-q01
Let the side be $1$ and drop the perpendicular $DP$ to $BC$. Since $\angle DBC=45^\circ$, $BP=DP$; since $\angle C=60^\circ$, $PC=\dfrac{DP}{\sqrt3}$.

\begin{asy}
usepackage("amsmath");
size(4.5cm);
defaultpen(fontsize(10pt));
path ra(pair P,pair u,pair v,real s=0.2){return P+s*unit(u)--P+s*unit(u)+s*unit(v)--P+s*unit(v);}
pair B=(0,0),C=(1,0),A=(0.5,sqrt(3)/2); real dp=sqrt(3)/(sqrt(3)+1); pair Dp=(dp,dp), P=(dp,0);
draw(A--B--C--cycle); draw(B--Dp,blue); draw(Dp--P,red+dashed); draw(ra(P,(1,0),(0,1),0.05));
label("$A$",A,N); label("$B$",B,SW); label("$C$",C,SE); label("$D$",Dp,NE); label("$P$",P,S); label("$45^\circ$",B+(0.13,0.05),fontsize(7));
label("\mbox{}",(min(currentpicture,user=true)+max(currentpicture,user=true))/2,black); // keeps the text after the figure black when LaTeX typesets the labels inline
\end{asy}

From $BP+PC=1$: $DP\left(1+\dfrac{1}{\sqrt3}\right)=1$, so $DP=\dfrac{\sqrt3}{\sqrt3+1}$. Then $[CDB]=\dfrac{1}{2}\cdot1\cdot DP$ and $[ABC]=\dfrac{\sqrt3}{4}$, so
\[
\dfrac{[ADB]}{[CDB]}=\dfrac{[ABC]}{[CDB]}-1=\dfrac{\sqrt3}{2DP}-1=\dfrac{\sqrt3+1}{2}-1=\boxed{\dfrac{\sqrt3-1}{2}}.
\]
